\documentclass{article}
\usepackage[portuguese]{babel}
\usepackage{listings}
\usepackage{xcolor}
\definecolor{codegreen}{rgb}{0,0.6,0}
\definecolor{codegray}{rgb}{0.5,0.5,0.5}
\definecolor{codepurple}{rgb}{0.8,0,0.2}
\definecolor{backcolour}{rgb}{.8,.8,1}
\lstdefinestyle{mystyle}{
    backgroundcolor=\color{backcolour},   
    commentstyle=\color{codegreen},
    keywordstyle=\color{blue},
    numberstyle=\tiny\color{codegray},
    stringstyle=\color{codepurple},
    basicstyle=\ttfamily\footnotesize,
    breakatwhitespace=false,         
    breaklines=true,                 
    captionpos=b,                    
    keepspaces=true,                 
    numbers=left,                    
    numbersep=5pt,                  
    showspaces=false,                
    showstringspaces=false,
    showtabs=false,                  
    tabsize=2
}
\lstset{style=mystyle}

\begin{document}
\section{Metodo da bisseccao}

	O metodo da bisseccao para encontrar raizes de funcoes, funciona dividindo pela metade o intervalo contendo a raiz ate que o intervalo seja tao pequeno que qualquer valor de $x$ pertencente a esse intervalo seja uma boa aproximacao da raiz verdadeira de acordo com a precisao escolhida.

    Comecamos com um intervalo inicial $[a,b]$, definimos um valor de $c$ no meio do intervalo$$c = \frac{a+b}{2}$$. Se $f(c)\times f(a)$ for negativo, sabemos que o zero esta entre $a$ e $c$ e entao trocamos b por $c$. Se o produto nao for negativo, sabemos que a raiz esta no intervalo $[c,b]$ e entao trocamos $a$ por $c$.Repetimos esses passos ate que alguma condicao de parada seja alcancada.
    
    O produto $f(a)\times f(b)$ ser negativo indica uma raiz no intervalo porque para que esse produto seja negativo um fator deve ser positivo e o outro negativo.


    
    
\begin{lstlisting}[language=Python, caption=Exemplo em python]
    def bissec(inf, sup, precisao, fun): 
    a = inf #definimos o limite inferior do intervalo inicial
    b = sup #limite superior
    x = 0   
    funcao = fun # f(x)

    while((b-a) < precisao): # testamos se o intervalo e suficienttemente pequeno
        x = (a+b)/2 # definimos um x no meio do intervalo
        y = funcao(x)

        if(y*funcao(a) < 0): # se f(x)*f(a) for negativo:
            b = x            #  trocamos b por x fazendo o intervalo [a,x]
        else:                # caso contrario: 
            a = x            #  trocamos a por x fazendo o intervalo [x,b]

        if(abs(y) < precisao): #testamos se f(x) esta proximo suficiente de 0
            break
    
    return x 
\end{lstlisting}

    Desde que o algoritmo base seja mantido, podemos mudar o codigo. Podemos a cada passo imprimir o valor de x, o intervalo, o numero de passos, definir uma condicao de parada depois de n passos, etc.

\begin{lstlisting}[language=C++, caption=Exemplo em C++]
typedef long double ddouble;

ddouble bissec(ddouble inf, ddouble sup, ddouble precisao, ddouble (*funcao)(ddouble){
    ddouble a = inf;
    ddouble b = sup;
    ddouble x;
    while(true){
        x = (a+b)/2;
        ddouble y = funcao(x);
        if((b-a)/2 < precisao){
            break;
        }if(abs(y) < precisao){
            break;
        }if

        if(funcao(a)*y < 0){
            b = x;
        }else{
            a = x;
        }

    };

    return(x);
};
\end{lstlisting}

Ainda que o codigo mude em diferentes linguagens, o algoritmo continua o mesmo.

\end{document}